% CLAUDE REPORT::BRIDGE-REFLECTION-0001::2026-11-04
\documentclass[11pt]{article}
\usepackage{amsmath,amssymb,amsthm}
\newtheorem{theorem}{Theorem}\newtheorem{corollary}{Corollary}
\title{CLAUDE BRIDGE-REFLECTION-0001: The weakest sufficient bridge $SAT\in P\Rightarrow \overline W_R\le \mathrm{poly}$ is exactly polynomial reflection of one SAT algorithm}
\date{November 4, 2026}
\begin{document}\maketitle
\noindent\textbf{Status.} Theorems 1--3: PROVED/EXACT (short, standard techniques; KNOWN-THEOREM flavour, no novelty claimed). Hard-instance lower bound for any strong $R$, and P versus NP: \textbf{OPEN}.

\section*{Setup}
$R$ is a fixed Cook--Reckhow proof system closed under substitution of constants and able to verify a polynomial-size circuit evaluation (Extended Frege, or the repository's $R_{\mathrm{mix341}}$ with DEF). $\overline W_R(N)=\max_{|F|\le N,\,F\in UNSAT}\min_{\pi}|\pi|$; the optimal $\pi$ is free. Bridge: $B3(R)$: $SAT\in P\Rightarrow \overline W_R(N)=N^{O(1)}$.

For a deterministic algorithm $A$ with time $n^c$ on inputs of length $n$, let $\mathrm{Sound}_{A,n}(f,\tau,y)$ be the propositional formula: ``if $\tau$ is a valid computation trace of $A$ on input bits $f$ ending in output UNSAT, then $f$ (parsed as CNF) is falsified by assignment $y$''. Its size is $\mathrm{poly}(n)$. It is a tautology iff $A$ never wrongly outputs UNSAT.

\begin{theorem}[Reflection equivalence]
Assume $A$ is a correct polynomial-time SAT decider. Then $\overline W_R(N)\le \mathrm{poly}(N)$ iff $\mathrm{Sound}_{A,n}$ has $R$-proofs of size $\mathrm{poly}(n)$. Hence $B3(R)\iff$ [$SAT\in P\Rightarrow$ for some correct poly-time $A$, $R$ has polynomial proofs of $\mathrm{Sound}_{A,n}$].
\end{theorem}
\begin{proof}
($\Leftarrow$) For UNSAT $F$ of length $n$ compute the trace $\tau_F$ of $A$; substitute the constants $f=F,\tau=\tau_F$ into the proof of $\mathrm{Sound}_{A,n}$. The antecedent becomes a constant circuit evaluation, refuted in $\mathrm{poly}(n)$ size by evaluation, leaving $\forall y\,\neg F(y)$, i.e.\ a refutation of $F$. ($\Rightarrow$) $\mathrm{Sound}_{A,n}$ is a tautology of size $\mathrm{poly}(n)$; $\overline W_R$ polynomial means $R$ is polynomially bounded on all tautologies (via negation), in particular on it.
\end{proof}

\begin{corollary}[Weak systems cannot carry the bridge]
If $\overline W_R$ is unconditionally superpolynomial (e.g.\ Resolution, Haken 1985), then $B3(R)\iff SAT\notin P$. If $B3(R)$ holds and $\overline W_R$ is superpolynomial then $SAT\notin P$. So a bridge is informative only for $R$ whose lower bound is itself open, and the pair (bridge, lower bound) must be proved for the \emph{same} $R$.
\end{corollary}

\begin{theorem}[Cost transfer to the repository's mixed system]
Since $R_{\mathrm{mix341}}$ contains charged DEF it p-simulates Extended Resolution (REFEREE299), and Extended Resolution p-simulates evaluation of circuits. Theorem 1 therefore applies to it with the same proof; the bridge for $R_{\mathrm{mix341}}$ reduces to a poly-size proof of one explicit tautology family $\mathrm{Sound}_{A,n}$ for \emph{some} correct $A$, and a lower bound for $\overline W_{R_{\mathrm{mix341}}}$ would give an Extended Resolution lower bound, a long-open problem.
\end{theorem}

\section*{Consequences for the programme}
(i) Searching hard UNSAT families against a strong $R$ is equivalent to showing $\mathrm{Sound}_{A,n}$ (for every candidate $A$) has no short $R$-proof: the hard instance is a \emph{reflection of a hypothetical algorithm}, which suggests restricting attention to $A$ with transparent structure (UWI portfolio, ARC/FLMV time--space algorithms) rather than generic $F$. (ii) Known sufficient route: if bounded arithmetic ($PV$ or $S^1_2$) proves $A$ correct then Extended Frege has the proofs (Cook; Krajíček--Pudlák); unprovability of P=NP-type statements in such theories is exactly the obstruction. (iii) Space/time audit: the trace has size $O(T_A(n)\cdot s_A(n))$ bits, so a time--space contraction of $A$ shrinks $|\mathrm{Sound}_{A,n}|$; this couples the ARC/FLMV spine to proof size but proves nothing about $R$-proof lower bounds.

\section*{Literature (this run)}
Cook--Reckhow 1979; Krajíček--Pudlák 1989; Haken 1985. Search found: Pich--Santhanam (arXiv 2312.08163, ``Towards P$\ne$NP from Extended Frege lower bounds'': split NP$=$coNP vs.\ $\ne$, relevant to our bridge case-split); Proof Analysis Problem (arXiv 2506.16956); Krajíček, arXiv 2604.04830 (feasible disjunction failure). No result gives $B3$ for a strong $R$. Classification: Pich--Santhanam = useful structural analogy requiring a transfer theorem; others = background.

\section*{Corpus audit}
Reviewed: AGENTS.md, LSPV-1/350 (MUS quotient), LSPV-4/386 (TABLE firewall, quantifier separation pointwise vs worst-case), claude ledger rows. Duplication: Theorem 1's $\Leftarrow$ direction overlaps REFLECTION336 (conditional); the new content is the \emph{equivalence} and Corollary 1. Stale claims: none found. Full main.tex ownership-region reconciliation and the other four roles' work were not performed in this run. Next arrow: exhibit or refute short $R_{\mathrm{mix341}}$-proofs of $\mathrm{Sound}_{A,n}$ for the UWI-portfolio $A$. P versus NP remains OPEN.
% END_OF_VERIFIED_REPORT::CLAUDE::BRIDGE-REFLECTION-0001
